\documentclass[aspectratio=169, xcolor={table,xcdraw}]{beamer}

\usepackage[T1]{fontenc}
\usepackage[utf8]{inputenc}
\usepackage{tgpagella}
\usepackage{tgcursor}
\usefonttheme{serif}
\usepackage[spanish, es-noshorthands]{babel}
\renewcommand{\tablename}{Tabla}

\usepackage{amsmath, amssymb, bm}
\usepackage{graphicx}
\usepackage{booktabs}
\usepackage[most]{tcolorbox}
\usepackage{tikz}
\usepackage{pgfplots}
\pgfplotsset{compat=1.18}
\usetikzlibrary{shapes.geometric, arrows.meta, positioning, calc}

\usetheme{Madrid}
\usecolortheme{default}
\definecolor{ucbblue}{RGB}{0, 51, 102}
\setbeamercolor{structure}{fg=ucbblue}
\setbeamercolor{title}{fg=white, bg=ucbblue}
\setbeamercolor{frametitle}{fg=white, bg=ucbblue}

\title[Singularidades y LSPB]{Transición a Singularidades y Generación de Trayectorias (LSPB)}
\subtitle{De la Velocidad Local al Movimiento Continuo}
\author[B. Quiroga Turdera]{M.Sc. Ing. Bernardo Quiroga Turdera}
\institute[UCB Tarija]{Universidad Católica Boliviana ``San Pablo'' — Sede Tarija}
\date{Semana 9 - Sesión 1}

\begin{document}

\begin{frame}
    \titlepage
\end{frame}

\begin{frame}{Diagnóstico: Lectura del Jacobiano[cite: 13]}
    \begin{columns}[T]
        \begin{column}{0.48\textwidth}
            \small
            Brazo 2R planar ($a_1=a_2=1$), evaluado en $q_1=0, q_2=\pi/2$. Jacobiano[cite: 13]:
            {\footnotesize
            \begin{equation*}
                J = \begin{bmatrix} -1 & -1 \\ 1 & 0 \\ 0 & 0 \\ 0 & 0 \\ 0 & 0 \\ 1 & 1 \end{bmatrix}, \quad J_p = \begin{bmatrix} -1 & -1 \\ 1 & 0 \end{bmatrix} \text{[cite: 13]}
            \end{equation*}
            }
            
            \textbf{Mapeo:} Si $\dot{\mathbf{q}} = [1, -1]^T$[cite: 13]:
            \begin{center}
                $\mathbf{v}_e = J\dot{\mathbf{q}} = [0 \;\; 1 \;\; 0 \;\; 0 \;\; 0 \;\; 0]^T$ \text{[cite: 13]}
            \end{center}
        \end{column}
        \begin{column}{0.5\textwidth}
            \begin{tcolorbox}[colback=white, colframe=ucbblue, title=Interpretación Física{[cite: 13]}, boxsep=0mm, top=1mm, bottom=1mm, left=1mm, right=1mm]
                \footnotesize
                \begin{itemize}
                    \item \textbf{Columna $J_i$:} Contribución de velocidad en el efector por velocidad unitaria en la junta $i$[cite: 13].
                    \item \textbf{Resultado $\mathbf{v}_e$:} $\mathbf{v}_L = [0,1,0]^T$, $\boldsymbol{\omega} = [0,0,0]^T$[cite: 13].
                    \item Movimientos articulares opuestos cancelan la velocidad angular $\implies$ \textbf{traslación instantánea pura} en $+y$[cite: 13].
                \end{itemize}
            \end{tcolorbox}
            \begin{alertblock}{No confundir}
                \footnotesize ¡La matriz $J$ NO es una transformación homogénea y $J\dot{\mathbf{q}}$ NO es una posición![cite: 13]
            \end{alertblock}
        \end{column}
    \end{columns}
\end{frame}

\begin{frame}{Rango Matricial y Movilidad Instantánea[cite: 13]}
    \begin{columns}[T]
        \begin{column}{0.48\textwidth}
            El \textbf{rango} mide el número de direcciones de velocidad independientes en el espacio de la tarea[cite: 13].
            
            Para $J \in \mathbb{R}^{m \times n}$, el rango máximo es[cite: 13]:
            \begin{equation*}
                \operatorname{rank}_{\max} = \min(m,n) \text{[cite: 13]}
            \end{equation*}
            
            \textbf{Ejemplo Rango Completo[cite: 13]:}
            \begin{equation*}
                J_A = \begin{bmatrix} -1 & -1 \\ 1 & 0 \end{bmatrix} \implies \operatorname{rank} J_A = 2
            \end{equation*}
        \end{column}
        \begin{column}{0.48\textwidth}
            \textbf{Ejemplo Rango Deficiente[cite: 13]:}
            \begin{equation*}
                J_B = \begin{bmatrix} 0 & 0 \\ 2 & 1 \end{bmatrix} \text{[cite: 13]}
            \end{equation*}
            Las columnas son proporcionales, por tanto[cite: 13]:
            \begin{equation*}
                \operatorname{rank} J_B = 1 \text{[cite: 13]}
            \end{equation*}
            
            \begin{tcolorbox}[colback=yellow!10, colframe=orange, title=Significado Físico, boxsep=0mm, top=1mm, bottom=1mm, left=1mm, right=1mm]
                \small La pérdida de rango \textbf{no} significa que el robot esté globalmente congelado, sino que se ha perdido una dirección de movimiento local independiente[cite: 13].
            \end{tcolorbox}
        \end{column}
    \end{columns}
\end{frame}

\begin{frame}{Singularidad Cinemática: Definición Universal vs. Especial[cite: 13]}
    \begin{tcolorbox}[colback=white, colframe=ucbblue, title=Definición Universal{[cite: 13]}, boxsep=0mm, top=1mm, bottom=1mm, left=1mm, right=1mm]
        Una singularidad cinemática es una configuración donde el Jacobiano pierde rango[cite: 13]:
        \begin{center}
            $\boxed{\operatorname{rank}J(\mathbf{q}) < \operatorname{rank}_{\max}}$ \text{[cite: 13]}
        \end{center}
    \end{tcolorbox}
    
    \vspace{0.3cm}
    \textbf{Caso Especial (Matrices Cuadradas)[cite: 13]:} \\
    \small Si $J$ es cuadrada (ej. $6 \times 6$ o el $J_p$ $2 \times 2$), una prueba equivalente es[cite: 13]:
    \begin{center}
        $\boxed{\det J(\mathbf{q}) = 0}$ \text{[cite: 13]}
    \end{center}
    
    \vspace{0.3cm}
    \begin{alertblock}{Precaución[cite: 13]}
        \footnotesize $\det J=0$ \textbf{no es una definición universal}. Para robots redundantes con Jacobianos no cuadrados (ej. $6 \times 7$), el determinante no existe, debiendo analizarse el rango directamente[cite: 13].
    \end{alertblock}
\end{frame}

\begin{frame}{Visualización de la Singularidad: Brazo Planar 2R[cite: 13]}
    \small El determinante del brazo 2R es $\det J_p = a_1 a_2 \sin q_2$[cite: 13]. Si $q_2 = 0$, $\det J_p = 0$.
    \vspace{0.2cm}
    
    \begin{columns}[T]
        \begin{column}{0.48\textwidth}
            \centering
            \textbf{Regular ($q_1=0, q_2=\pi/2$)[cite: 13]}
            \vspace{0.1cm}
            
            \begin{tikzpicture}[scale=0.85, thick]
                \draw[->, gray] (0,0) -- (1.5,0);
                \draw[->, gray] (0,0) -- (0,1.5);
                
                \draw[ucbblue, line width=2pt] (0,0) -- (1,0) -- (1,1);
                \filldraw[black] (0,0) circle (1pt); \filldraw[black] (1,0) circle (1pt); \filldraw[red] (1,1) circle (1pt);
                
                \draw[-Stealth, red!70!black, line width=1.5pt] (1,1) -- ++(-0.5, 0.5) node[above]{$J_{1}$};
                \draw[-Stealth, orange, line width=1.5pt] (1,1) -- ++(-0.5, 0) node[left]{$J_{2}$};
            \end{tikzpicture}
            
            \footnotesize Columnas independientes.
        \end{column}
        \begin{column}{0.48\textwidth}
            \centering
            \textbf{Singular ($q_1=0, q_2=0$)[cite: 13]}
            \vspace{0.1cm}
            
            \begin{tikzpicture}[scale=0.85, thick]
                \draw[->, gray] (0,0) -- (2.5,0);
                \draw[->, gray] (0,0) -- (0,1);
                
                \draw[ucbblue, line width=2pt] (0,0) -- (1,0) -- (2,0);
                \filldraw[black] (0,0) circle (1pt); \filldraw[black] (1,0) circle (1pt); \filldraw[red] (2,0) circle (1pt);
                
                \draw[-Stealth, red!70!black, line width=1.5pt] (2,0) -- ++(0, 1) node[right]{$J_{1}$};
                \draw[-Stealth, orange, line width=1.5pt] (2,0) -- ++(0, 0.5) node[left]{$J_{2}$};
            \end{tikzpicture}
            
            \footnotesize Columnas paralelas (en $+y$). \\
            \textbf{Pérdida de dirección radial ($x$)[cite: 13]}.
            
        \end{column}
    \end{columns}
    
    \vspace{0.2cm}
    \footnotesize \textit{Conexión IK Global: La misma postura marca la frontera del espacio de trabajo (no puede alcanzar $x > a_1+a_2$)[cite: 13].}
\end{frame}

\begin{frame}{Cinemática Inversa Diferencial cerca de Singularidades[cite: 13]}
    \begin{columns}[T]
        \begin{column}{0.48\textwidth}
            Para un Jacobiano cuadrado no singular, la IK diferencial es[cite: 13]:
            \begin{equation*}
                \boxed{\dot{\mathbf{q}} = J^{-1}(\mathbf{q})\dot{\mathbf{x}}_d} \text{[cite: 13]}
            \end{equation*}
            
            ¿Qué ocurre al aproximarse a una singularidad? Sea el modelo ilustrativo[cite: 13]:
            \begin{equation*}
                J_\varepsilon = \begin{bmatrix} 1 & 0 \\ 0 & \varepsilon \end{bmatrix} \implies J_\varepsilon^{-1} = \begin{bmatrix} 1 & 0 \\ 0 & 1/\varepsilon \end{bmatrix} \text{[cite: 13]}
            \end{equation*}
        \end{column}
        \begin{column}{0.48\textwidth}
            Si deseamos una velocidad unitaria $\dot{\mathbf{x}}_d = [0, 1]^T$[cite: 13]:
            \begin{equation*}
                \dot{\mathbf{q}} = \begin{bmatrix} 0 \\ 1/\varepsilon \end{bmatrix} \text{[cite: 13]}
            \end{equation*}
            
            \begin{center}
                \small
                \begin{tabular}{r | l}
                $\varepsilon$ & Velocidad requerida $\dot{q}_2$ \\ \hline
                1 & 1 \\
                0.1 & 10 \\
                0.01 & 100
                \end{tabular}
            \end{center}
            \begin{tcolorbox}[colback=yellow!10, colframe=orange, title=Significado Físico, boxsep=0mm, top=1mm, bottom=1mm, left=1mm, right=1mm]
                \small Conclusión: Cerca de la singularidad, mantener una velocidad cartesiana finita demanda velocidades articulares inalcanzables[cite: 13].
            \end{tcolorbox}
        \end{column}
    \end{columns}
\end{frame}

\begin{frame}{De Puntos a Trayectorias (Planificación)[cite: 13]}
    Hasta ahora la IK nos daba la configuración de destino. Pero dos puntos ($q_A, q_B$) no definen una ley de movimiento[cite: 13].
    
    \vfill
    \begin{center}
    \begin{tikzpicture}[scale=1, thick]
        \node[draw, rounded corners, fill=ucbblue!10, text width=3cm, align=center] (C) at (0,0) {\textbf{Configuración}\\Solo los puntos extremos.};
        \node[draw, rounded corners, fill=ucbblue!10, text width=3cm, align=center] (P) at (4.5,0) {\textbf{Camino (Path)}\\Curva geométrica sin temporización.};
        \node[draw, rounded corners, fill=green!10, text width=4cm, align=center] (T) at (9.5,0) {\textbf{Trayectoria}\\Parametrizada en el tiempo: $q(t), \dot{q}(t), \ddot{q}(t)$.};
        
        \draw[->] (C) -- (P);
        \draw[->] (P) -- (T);
    \end{tikzpicture}
    \end{center}
    
    \vfill
    \textbf{¿Por qué es vital el tiempo?}[cite: 13] \\
    Los extremos por sí solos no determinan la velocidad, la aceleración requerida ni, en última instancia, el esfuerzo del motor ($\tau$)[cite: 13].
\end{frame}

\begin{frame}{Generación de Trayectorias: Perfil LSPB (1/2)[cite: 13]}
    \small 
    \textbf{LSPB} (Linear Segment with Parabolic Blends) genera un perfil de velocidad \textbf{trapezoidal} para un movimiento punto a punto de una junta[cite: 13].
    
    \textbf{Suposiciones y Fases[cite: 13]:} Posiciones $q_i, q_f$; tiempo final $t_f$; reposo en extremos $\dot{q}(0)=\dot{q}(t_f)=0$; aceleración constante $\ddot{q}_c$[cite: 13].
    
    \vspace{0.4cm}
    \begin{columns}[T]
        \begin{column}{0.33\textwidth}
            \footnotesize
            \textbf{1. Aceleración} ($0 \le t \le t_c$)\\
            $q(t) = q_i + \tfrac{1}{2}\ddot{q}_c t^2$ \\
            $\dot{q}(t) = \ddot{q}_c t$ \\
            $\ddot{q}(t) = \ddot{q}_c$
        \end{column}
        \begin{column}{0.34\textwidth}
            \footnotesize
            \textbf{2. V. Constante} ($t_c < t \le t_f - t_c$)\\
            $q(t) = q_i + \dot{q}_c (t - \tfrac{t_c}{2})$ \\
            $\dot{q}(t) = \ddot{q}_c t_c$ \\
            $\ddot{q}(t) = 0$
        \end{column}
        \begin{column}{0.33\textwidth}
            \footnotesize
            \textbf{3. Desaceleración} ($t_f - t_c < t \le t_f$)\\
            $q(t) = q_f - \tfrac{1}{2}\ddot{q}_c(t_f - t)^2$ \\
            $\dot{q}(t) = \ddot{q}_c(t_f - t)$ \\
            $\ddot{q}(t) = -\ddot{q}_c$
        \end{column}
    \end{columns}
\end{frame}

\begin{frame}{Generación de Trayectorias: Perfil LSPB (2/2)[cite: 13]}
    \textbf{Relación de Tiempo de Mezcla ($t_c$)[cite: 13]:}
    \begin{center}
        $\boxed{ t_c = \frac{t_f}{2} - \frac{1}{2}\sqrt{t_f^2 - \frac{4(q_f - q_i)}{\ddot{q}_c}} }$ \text{[cite: 13]}
    \end{center}
    
    \vfill
    \textbf{Restricción Física (Aceleración Mínima)[cite: 13]:}
    Para que el movimiento sea factible en el tiempo estipulado, el motor debe cumplir[cite: 13]:
    \begin{center}
        $\boxed{ |\ddot{q}_c| \ge \frac{4|q_f - q_i|}{t_f^2} }$ \text{[cite: 13]}
    \end{center}
    \vfill
    \textit{Si se usa la igualdad, el segmento de velocidad constante desaparece y el perfil de velocidad se vuelve puramente triangular[cite: 13].}
\end{frame}

\begin{frame}{Ejemplo Guiado: LSPB[cite: 13]}
    \begin{columns}[T]
        \begin{column}{0.48\textwidth}
            \textbf{Datos[cite: 13]:}
            $q_i=0$, $q_f=2$ rad, $t_f=3$ s, $\ddot{q}_c=1$ rad/s$^2$.
            
            \vspace{0.2cm}
            \textbf{1. Viabilidad[cite: 13]:}
            $|\ddot{q}_c|_{\min} = \frac{4(2)}{3^2} = \frac{8}{9}$. Como $1 > 8/9$, es factible.
            
            \vspace{0.2cm}
            \textbf{2. Tiempos y Velocidad[cite: 13]:}
            $t_c = \frac{3}{2} - \frac{1}{2}\sqrt{9-8} = 1$ s. \\
            $\dot{q}_c = \ddot{q}_c t_c = 1$ rad/s.
            
            \vspace{0.2cm}
            \textbf{3. Ecuaciones por tramos[cite: 13]:}
            \begin{align*}
                \dot{q}(t) &= \begin{cases} t & 0 \le t \le 1 \\ 1 & 1 < t \le 2 \\ 3-t & 2 < t \le 3 \end{cases}
            \end{align*}
        \end{column}
        \begin{column}{0.5\textwidth}
            \centering
            \begin{tikzpicture}[scale=0.85]
                \begin{axis}[
                    width=7cm, height=4.5cm,
                    xlabel={$t$ [s]}, ylabel={$\dot{q}(t)$ [rad/s]},
                    xmin=0, xmax=3, ymin=0, ymax=1.2,
                    xtick={0,1,2,3}, ytick={0,1},
                    grid=both,
                    title={Perfil de Velocidad Trapezoidal}
                ]
                \addplot[ucbblue, thick, domain=0:1] {x};
                \addplot[ucbblue, thick, domain=1:2] {1};
                \addplot[ucbblue, thick, domain=2:3] {3-x};
                
                \draw[dashed, red] (axis cs:1,0) -- (axis cs:1,1);
                \draw[dashed, red] (axis cs:2,0) -- (axis cs:2,1);
                \end{axis}
            \end{tikzpicture}
            
            \footnotesize $q(t)$ y $\dot{q}(t)$ son continuas. $\ddot{q}(t)$ presenta saltos en los límites de fase[cite: 13].
        \end{column}
    \end{columns}
\end{frame}

\begin{frame}{Arquitectura Integrada del Hito 2 (Pipeline)[cite: 13]}
    \begin{center}
    \begin{tikzpicture}[
        node distance=0.4cm and 0.5cm,
        box/.style={rectangle, draw=ucbblue, thick, fill=ucbblue!5, text width=2.4cm, align=center, rounded corners, font=\scriptsize},
        arrow/.style={-Stealth, thick, ucbblue}
    ]
        \node[box] (ik) {\textbf{Pieper (IK)}\\$T_d \to \mathbf{q}$};
        \node[box, right=of ik] (fk) {\textbf{Validación}\\Selección + FK};
        \node[box, right=of fk] (lspb) {\textbf{LSPB}\\$q(t), \dot{q}(t), \ddot{q}(t)$};
        \node[box, below=of lspb] (jac) {\textbf{Jacobiano}\\$J(q(t))$};
        \node[box, left=of jac] (sing) {\textbf{Singularidades}\\Evaluación Local};
        \node[box, fill=yellow!20, left=of sing] (dyn) {\textbf{Dinámica}\\(Pendiente W09-S02)};
        \node[box, fill=green!10, below=of dyn] (tau) {\textbf{Torque}\\$\tau(t)$};
        
        \draw[arrow] (ik) -- (fk);
        \draw[arrow] (fk) -- (lspb);
        \draw[arrow] (lspb) -- (jac);
        \draw[arrow] (jac) -- (sing);
        \draw[arrow] (sing) -- (dyn);
        \draw[arrow] (dyn) -- (tau);
    \end{tikzpicture}
    \end{center}
    
    \vspace{0.2cm}
    \textbf{Preparación para W09-S02 (Dinámica)[cite: 13]:} \\
    W09-S01 nos entrega $q(t), \dot{q}(t), \ddot{q}(t)$. La próxima sesión ingresará esos datos a[cite: 13]:
    \begin{center}
        $\boxed{\tau = M(q)\ddot{q} + C(q, \dot{q})\dot{q} + g(q)}$ \text{[cite: 13]}
    \end{center}
\end{frame}

\end{document}
